%!TEX TS-program = xelatex
\documentclass[]{friggeri-cv}
\usepackage[brazil,english]{babel}
\usepackage{xltxtra}
\defaultfontfeatures{Ligatures=TeX}
\usepackage{afterpage}
\usepackage{hyperref}
\usepackage{hyphenat}
\usepackage{color}
\usepackage{xcolor}
\usepackage{fancyhdr}
\usepackage{multilanguage}
\usepackage{marvosym}  % used for glyphs such as email, mobile phone, etc.

%\setdoclang{br}{brazil}
\setdoclang{en}{english}

\hypersetup{%
    colorlinks,%
    linkcolor={red!50!black},%
    citecolor={blue!50!black},%
    urlcolor={blue!80!black},%
    pdfborder={0 0 0},%
}

\pagestyle{fancy}
\lhead{}
\chead{}
\rhead{}
\rfoot{{%
    \footnotesize{%
        \langif{br}{%
        %lang pt-br
            Veja a última versão deste documento em:
            \href{https://github.com/diraol/cv/releases/}{https://github.com/diraol/cv/releases/}%
        }{%
        %lang en
            See the lastest version of this document at:
            \href{https://github.com/diraol/cv/releases/}{https://github.com/diraol/cv/releases/}%
        }
}}}
\renewcommand{\headrulewidth}{0pt}
\renewcommand{\footrulewidth}{0pt}

\hypersetup{%
    pdftitle={Diego Rabatone Oliveira Resume},%
    pdfauthor={Diego Rabatone Oliveira},%
    pdfsubject={Resume/CV},%
    pdfkeywords={Diego Rabatone Oliveirar, diraol, cv, resume},%
    colorlinks=false,       % no lik border color
    allbordercolors=white    % white border color for all
}
%\addbibresource{bibliography.bib}
\RequirePackage{xcolor}
\definecolor{pblue}{HTML}{0395DE}

\begin{document}
\thispagestyle{empty}
\pagenumbering{gobble}% Remove page numbers (and reset to 1)
%%%%%%%%%%%%%%%%%%%%%%%%%%%%%%%%%%%%%%%%%%
%\langif{br}{%
%% text for lang pt-br
%
%}{%
%% text for lang pt-br
%
%}
%\langif{br}{}{}
%%%%%%%%%%%%%%%%%%%%%%%%%%%%%%%%%%%%%%%%%%
\header{Diego}{Rabatone O.}{\langif{br}{Engenheiro}{Engineer}}

% Fake text to add separator
\fcolorbox{white}{gray}{\parbox{\dimexpr\textwidth-2\fboxsep-2\fboxrule}{%

}}
%
% LATERAL
%
% In the aside, each new line forces a line break
\begin{asidefp}
  \section{\langif{br}{Contato}{Contact}}
    \pin \hfill São Paulo, SP, \langif{br}{Brasil}{Brazil}%
~
    {\Large\textcolor{gray}{\Mobilefone}} \hfill +55 11 9-8231-4249%
~
    {\Large\textcolor{gray}{\Letter}} \hfill \href{mailto:cv@diraol.eng.br}{cv@diraol.eng.br}%
    % {\Large\textcolor{gray}{\Letter}} \hfill \href{mailto:diraol@alumni.usp.br}{diraol@alumni.usp.br}%
~
    \tlogo \hfill \href{http://twitter.com/diraol}{@diraol}%
~
    \llogo \hfill \langif{br}{\href{https://br.linkedin.com/in/diraol}{in://diraol}}{\href{https://www.linkedin.com/in/diraol/}{in://diraol}}%
~
  %\section{IRC}
  %  \#diraol (at) irc.freenode.net
  %  ~
  \section{Web \& Git}
    % \href{http://cv.diraol.eng.br}{cv.diraol.eng.br}
    \href{https://gitlab.com/diraol}{gitlab.com/diraol}
    \href{https://github.com/diraol}{github.com/diraol}
    ~%
  \section{\langif{br}{Idiomas}{Languages}}
    \textbf{\langif{br}{Português}{Portuguese}}\includegraphics[scale=0.40]{img/5stars.png}
    \textbf{\langif{br}{Inglês}{English}}\includegraphics[scale=0.40]{img/5stars.png}
    \textbf{\langif{br}{Espanhol}{Spanish}}\includegraphics[scale=0.40]{img/2stars.png}
    \textbf{\langif{br}{Francês}{French}}\includegraphics[scale=0.40]{img/1stars.png}
\end{asidefp}
%
% EDUCAÇÃO
%
\sectionlang{br}{Educação}
\sectionlang{en}{Education}
\begin{entrylist}
  \entry
    {2005 - 2014}
    {\langif{br}{Graduação em Engenharia Elétrica}{Undergraduate in Electrical Engineering}}
    {Universidade de São Paulo, Brazil}
    {\langif{br}{Ênfase Computação - Concluída em 2014}{Minor in Computer and Digital Systems - Concluded in 2014}%
    %Main subjects: Mathematics and Physics, Programming, Operational Research, Telecommunication Systems, Digital and Analogical Electronics.\\
    %\emph{Title of the Thesis: "Development, Management and Migrations of web contents and applications".}\\
    %\emph{Thesis activity carried out during an internship period at Atitlan Engineering SRL.}\\
    }%
  %\entry
  %  {2000 - 2002}
  %  {\langif{br}{Ensino Médio}{High School}}
  %  {Colégio Rio Branco, São Paulo, Brasil}
  %  {%Scientific Secondary School.\\
  %  %Main subjects: Mathematics, Physics, Computer Science.
  %  }%
\end{entrylist}%
%
% EXPERIÊNCIA (parte 1)
%
\sectionlang{br}{Experiência}
\sectionlang{en}{Experience}
\begin{entrylist}
  \entry
    {\langif{br}{Dez/19 - Now}{Dec/19 - Now}}
    {DevOps Consultant}
    {\href{http://kobold.com.br/}{KOBOLD}}
    {\langif{br}{%
         Consultor SRE/DevOps com foco em padronização e automação de ambiente
         de desenvolvimento (CI/CD) e infraestrutura na AWS. Desenvolvimento de
         soluções para redução de custos e treinamento de equipe de
         desenvolvimento para uso \emph{self-service} de infraestrutura.%
       }{%
         Kobold is a Financial institution that provide a type of fund composed
         of receiveables. At Kobold I play a SRE/DevOps consultant role with
         focus on automation and standardization of development environment
         (CI/CD) and infrastructure on AWS. Cost reduction solutions and
         engineering training (both technical and methodological) to
         \emph{self-service} on the whole pipeline.%
       } %
     \maintechs
     AWS, Jenkins, Github, Packer, Terraform, Docker, shellscript.
     \hfill%
     \href{https://kobold.com.br/}{https://kobold.com.br/}%
    }
  \entry
    {\langif{br}{Ago/15 - Hoje}{Aug/15 - now}}
    {\langif{br}{Co-fundador}{Co-Founder}}
    {\href{https://ask-ar.xyz}{ASK-AR}}
    {\langif{br}{%
         Consultoria de análise de dados e desenvolvimento de tecnologia para
         promover maior equidade social, com excelência, crítica e ética.
       }{%
         ASK-AR is a consultancy company that delivers knowledge using data and
         tecnology working together with our clients. As ASK-AR founder I play a
         more strategical role on projects conception and planning.
       }%
     \hfill%
     \href{https://ask-ar.xyz}{https://ask-ar.xyz}%
    }
  \entry
    {\langif{br}{Dez/17 - Out/21}{Dec/17 - Oct/21}}
    {Senior Infrastructure Engineer}
    {\href{http://www.cabify.com/}{Cabify} / EasyTaxi}
    {\langif{br}{%
         Planejamento, desenvolvimento e manutenção da infraestrutura de cloud
         com foco em habilitar o time de engenharia para desenvolver e realizar
         deploy da plataforma de microserviços da Cabify. Migração da plataforma
         do GCP para AWS, gestão de instância própria do Gitlab com completa
         automação de pipelines de CI/CD, Prometheus, Cortex e Grafana como
         plataforma de observabilidade e clusters d Elasticsearch para logs.
         Infraestrutura gerenciada com uso de pipelines de Terraform, Packer e
         K8s (GKE e EKS). Membro do grupo de segurança.%
       }{%
         Planning, development and maintenance of Cloud Infrastructure to enable
         engineers to develop and deploy Cabify's microservice platform.
         Migration of the platform from GCP to AWS (2020/2021), management of
         selfhosted Gitlab with full CI/CD automated pipelines, Prometheus,
         Cortex and Grafana observability platform, Elasticsearch clusters for
         logs. Pipeline based Terraform managed infrastructure and K8s (GKE and
         EKS). Member of the Security guild.%
       } %
     \maintechs
     K8s, Packer, Terraform, Prometheus, Grafana, Elasticsearch, Gitlab, AWS \eand
     GCP.
     \hfill%
     \href{http://www.cabify.com/}{http://www.cabify.com/}%
    }%
  \entry%
    {}
    {\langif{br}{Engenheiro de Dados Sênior}{Senior Data Engineer}}
    {}
    {\langif{br}{%
         Responsável pelo time de Data Engineering da EasyTaxi. Gestão e
         evolução de todo pipeline de dados, Data Lake e Data Wareouse.
         Interlocução com stakeholders internos como Cientistas de dados,
         Analistas de Negócios, equipes de desenvolvimento, equipes de
         operações, dentre outros.
         Atuei também no projeto de integração das plataformas EasyTaxi e
         Cabify.%
       }{%
         Responsible for EasyTaxi Data Engineering team. Data pipelines, Data
         Lake and Data Warehouse. Close contact with many internal stakeholders,
         such as Data Scientists, Business Analysts, Development teams,
         Operational teams, among others.  Also worked with the integration
         between EasyTaxi and Cabify integration.%
       } \\ %
     \maintechs
     Kafka, Airflow, Python, Databricks, AWS, BigQuery, \eand InfluxDB.
     %\hfill%
     %\href{http://www.easytaxi.com/}{http://www.easytaxi.com/}%
    }
  % https://indico.cern.ch/event/647995/contributions/2652502/attachments/1507421/2349378/RENAFAE_2017_RLIope.pdf
  % https://web.archive.org/web/20190505190330/https://indico.cern.ch/event/647995/contributions/2652502/attachments/1507421/2349378/RENAFAE_2017_RLIope.pdf
  % https://sc17.supercomputing.org/
  % https://www.amlight.net/?p=2999
  % https://web.archive.org/web/20190505183456/https://www.amlight.net/?p=2999
  % https://web.archive.org/web/20190203153054/https://tnc17.geant.org/
  % https://indico.cern.ch/event/497415/contributions/1176934/attachments/1261884/1865776/INFIERI_Lisbon_RLIope.pdf
  % https://web.archive.org/web/20190505190602/https://indico.cern.ch/event/497415/contributions/1176934/attachments/1261884/1865776/INFIERI_Lisbon_RLIope.pdf
  \entry
    {\langif{br}{Mar/16 - Hoje}{Mar/16 - Nov/17}}
    {\langif{br}{Pesquisador}{Research Fellow}}
    {\href{http://ncc.unesp.br}{NCC-UNESP}}
    {\langif{br}{%
         Pesquisa e desenvolvimento em Redes Definidas por Software (SDN).\\
         Participação no time principal do \href{https://kytos.io}{Projeto Kytos}, atuando diretamente na concepção da arquitetura do projeto e em seu desenvolvimento.\\
         Palavras Chave:
       }{%
         Software-Defined Networking (SDN) R\&D.
         \href{https://kytos.io}{Kytos Project} core team member, working on the project architectural design and development,
         including the python-openflow library and other network applications from Kytos Ecosystem.\\
       }%
     \maintechs
     Python, SDN, OpenFlow, REST, API \eand d3js%
     \hfill%
     \href{https://kytos.io}{https://kytos.io}%
    }
\end{entrylist}
%
% NOVA PÁGINA
%
\newpage
\pagenumbering{gobble}% Remove page numbers (and reset to 1)
%
% BARRA LATERAL
%
\begin{aside}
~
~
~
  \section{\langif{br}{Tecnologias e ferramentas}{Technologies and tools}}
    awk
    AWS
    Airflow
    BigQuery
    Databricks
    Docker
    GCP
    Git
    Gitlab
    grep
    jq
    K8s
    Kafka
    Packer
    Python
    R
    sed
    Shell Script
    Terraform
    % CSS
    % JavaScript
    %\LaTeX
    %CartoDB
    % Django
    % Drupal
    % Gimp
    % Inkscape
    % Mediawiki
    % MySQL/MariaDB
    % NFS
    %Nginx
    % OpenLDAP
    % OpenFlow
    % Postgres
    % QuantumGis
    % SDN
    % Varnish
    % Wordpress
    % Xen
  ~
  \section{\langif{br}{Ambiente diário}{Daily Env}}
    Debian GNU/Linux \\%\includegraphics[scale=0.40]{img/5stars.png}
    swaywm \\%Gnome3 \\%\includegraphics[scale=0.40]{img/2stars.png}
    Vim
  ~
  \section{Hobbies}
      \href{http://olhares.com/diraol}{\langif{br}{Fotografia}{Photography}}
      Tae-Kwon-Do
  ~
  \section{\langif{br}{Cursos e Certificações}{Courses and Certifications}}
      \langif{br}%
      {\textbf{Especialização em Data Science}\\{\scriptsize{Universidade Johns Hopkins no Coursera}}}%
      {\textbf{Data Science Specialization}\\{\scriptsize{by Johns Hopkins University on Coursera}}}
      ~
      \langif{br}%
      {\textbf{Objetos Pythonicos}\\{\scriptsize{por \href{https://www.python.pro.br/modulos/}{PythonPro}}}}%
      {\textbf{Pythonic Objects}\\{\scriptsize{by \href{https://www.python.pro.br/modulos/}{PythonPro}}}}
\end{aside}
%
% EXPERIÊNCIA (parte 2)
%
\sectionlang{br}{Experiência - cont.}
\sectionlang{en}{Experience - cont.}
\begin{entrylist}
  \entry
    {\langif{br}{Out/09 - Fev/17}{Oct/09 - Feb/17}}
    {\langif{br}{Co-fundador, coordenador e membro}{Co-Founder, coordinator and member}}
    {\href{http://polignu.org}{PoliGNU}}
    {\langif{br}{Grupo de Estudos de Software Livre da Poli-USP}{Poli-USP Free Software Studies Group} (\href{https://polignu.org}{PoliGNU})\\
     \langif{br}{%
         Organização e planejamento, Gestão de Infraestrutura (SysAdmin e website).
         Apresentações em conferências, seminários, workshops, etc.\\
         Palavras Chave: Software Livre, projetos, divulgação%
       }{%
         Organization and planning, infrastructure management (SysAdmin, website),
         disclosure and presentation at conferences, seminars, workshops, etc.\\
         Keywords: Free Software, projects, advocating%
       }
     \hfill
     \href{http://polignu.org}{http://polignu.org}%
    }
  \entry
    {\langif{br}{Nov/15 - Fev/16}{Nov/15 - Feb/16}}
    {\langif{br}{Desenvolvedor}{Developer}}
    {\href{https://okfn.org/}{OKF}}
    {\langif{br}{
         Desenvolvimento de aplicação web com foco em Dados, em parceria
         com uma bolsista da Open Knowledge Foundation (OKF) das Filipinas.\\
       }{
         Development of data web application with an
         Open Knowledge Foundation (OKF) fellow from Philippines.\\
       }%
     \maintechs
     NodeJS, Heroku \eand PostgreSQL.
     %\href{http://washresilience.org/}{http://washresilience.org/}
    }
  \entry
    {\langif{br}{Ago/15 - Out/15}{Aug/15 - Oct/15}}
    {\langif{br}{Pesquisador Sênior }{Senior Researcher}}
    {\href{https://fga.unb.br/lappis}{LAPPIS/UNB e CGU}}
    {\langif{br}{%
         Especificação de aplicação para gestão de transparência governamental.\\
         Palavras Chave: Python/Django, Dados Governamentais Abertos%
       }{%
         Designed Software to control governamental transparency indexes.\\
         Keywords: Python, Django, OpenData, OpenGovernment, Transparency%
       }\\
     \href{https://gitlab.com/mibt/mibt}{https://gitlab.com/mibt/mibt}
   }
  \entry
    {\langif{br}{Ago/15 - Out/15}{Aug/15 - Oct/15}}
    {\langif{br}{Consultor}{Consultant}}
    {\href{http://www.gvces.com.br/}{FGV - GVCes}}
    {\langif{br}{%
         Atuação como \textit{Product Owner} no desenvolvimento de um sistema de Gestão e Divulgação de Indicadores
         ligado ao projeto \href{http://indicadoresdebelomonte.com.br}{``Indicadores de Belo Monte''}.
       }{%
         \textit{Product Owner} on the open data platform%
         \href{http://indicadoresdebelomonte.com.br}{``Indicadores de Belo Monte''}.
       }\\
     \href{http://indicadoresdebelomonte.com.br}{http://indicadoresdebelomonte.com.br}
     }
  \entry
    {\langif{br}{Fev/15 - Mai/15}{Feb/15 - May/15}}
    {\langif{br}{Consultor PNUD}{UNDP Consultant}}
    {\href{http://participacao.mj.gov.br}{\langif{br}{Ministério da Justiça}{Brazilian Ministry of Justice}}}
    {\langif{br}{Consultor pelo Programa das Nações Unidas para o Desenvolvimento}{United Nations Development Programme IT Consultant}\\
     \langif{br}{%
         Planejamento e desenvolvimento das plataformas de Consulta Pública do Ministério da Justiça.
         Implementação de metodologias ágeis.\\
         Palavras Chave: Wordpress, PHP, democracia participativa, métodos ágeis
       }{%
         Planning and development of Brazilian Justice Ministry Public Feedback Platforms.
         Implementation of agile methodology.\\
         Keywords: Wordpress, PHP, participative democracy, Agile
       }\\
     \href{http://pensandoodireito.mj.gov.br/}{http://pensandoodireito.mj.gov.br/} \& \href{http://github.com/pensandoodireito}{http://github.com/pensandoodireito}
    }
  \entry
    {\langif{br}{Jul/12 - Fev/15}{Jul/12 - Feb/15}}
    {\langif{br}{Analista de Mídias Digitais}{Digital Media Analyst}}
    {O Estado de S.Paulo S/A}
    {\langif{br}{%
         Desenvolvedor na equipe de jornalismo de dados.%
       }{%
         Main developer of the data driven journalism team.%
       } ("\href{http://estadaodados.com}{Estadão Dados}")\\
     \langif{br}{%
         Palavras Chave: extração, tratamento, união, processamento, análise e \nohyphens{visualização} de dados.%
       }{%
         Keywords: Data extraction, scraping, gathering, processing, analysis and \nohyphens{visualization}.%
       }\hfill
     \href{http://estadaodados.com}{http://estadaodados.com}
    }
   %\entry
   % {\langif{br}{Jul/08 - Dez/13}{Jul/08 - Dec/13}}
   % {\langif{br}{Administrador de Sistemas voluntário}{Volunteer SysAdmin}}
   % {Escritório Piloto - Poli-USP}
   % {\langif{br}{
   %      Implementação e manutenção de rede local (LAN) e website.\\
   %      Palavras Chave:
   %    }{
   %      Design and Setup of local network (LAN) and website.\\
   %      Keywords:
   %    }
   %  OpenLDAP, NFS, apache, XEN, DNS , Drupal.
   % }
   %%\entry
   %% {\langif{br}{Ago/10 - Ago/11}{Aug/10 - Aug/11}}
   %% {\langif{br}{Professor de Matemática}{Math Teacher}}
   %% {ACEPUSP}
   %% {\langif{br}{
   %%      Curso preparatório para vestibular para \nohyphens{estudantes} de baixa renda.\\
   %%      Palavras Chave: Matemática, Geometria, Trigonometria.
   %%    }{
   %%      University Preparatory course for low income students.\\
   %%      Keywords: Mathematics, Geometry, Trigonometry.
   %%    }
   %%  }
   %\entry
   % {\langif{br}{Ago/10 - Dez/10}{Aug/10 - Dec/10}}
   % {\langif{br}{Voluntário}{Volunteer}}
   % {Universidade de São Paulo}
   % {\langif{br}{Modelagem de base de dados}%para uma nova versão do Sistema de \nohyphens{Avaliação} de Ensino}%
   % {Database modeling}\\% for the new version of the Teaching Evaluation System}\\%
   % \href{https://github.com/SuperNovaPOLIUSP}{https://github.com/SuperNovaPOLIUSP}}
   %\entry
   % {\langif{br}{Jan/09 - Ago/10}{Jan/09 - Aug/10}}
   % {\langif{br}{Voluntário}{Volunteer}}
   % {Escola de Governo}
   % {\langif{br}{Desenvolvimento de website institucional (Joomla)}{Design, development and management of institucional website on Joomla}\\
   %  \langif{br}{Orientação e mentoria de estudantes do curso de Formação de Governantes}{Orientation and Mentoring of Governor Formation Course Students}}
   %\entry
   % {\langif{br}{Mar/05 - Jun/06}{Mar/05 - Jun/06}}
   % {\langif{br}{Bolsista}{Fellow}}
   % {Universidade de São Paulo}
   % {\langif{br}{
   %     Desenvolvedor do primeiro Sistema de Avaliação de Ensino da Escola \nohyphens{Politécnica}.
   %   }{
   %     Developer of the first Digital Teaching Evaluation System from Escola \nohyphens{Politécnica} da USP.}}
\end{entrylist}%
%
% PROJETOS
%
%\sectionlang{br}{Projetos e Comunidades}
%\sectionlang{en}{Projects and Communities}
\sectionlang{br}{Projetos}
\sectionlang{en}{Projects}
\textbf{python-mode} \hfill \small{\href{https://github.com/python-mode/python-mode}{https://github.com/python-mode/python-mode/}}\\
\small{\langif{br}%
{Mantenedor do plugin \emph{python-mode} para vim desde 2017.}%
{Maintainer of the vim plugin \emph{python-mode} since 2017.}}\\

\textbf{\langif{br}{Radar Parlamentar}{The Parliamentary Radar}} \hfill \small{\href{https://radarparlamentar.polignu.org/}{https://radarparlamentar.polignu.org/}}\\
\small{\langif{br}%
{Aplicativo que determina ``semelhanças'' entre partidos \nohyphens{políticos} baseado na análise matemática dos dados de votações de projetos de lei na casa \nohyphens{legislativa}. Existem ainda algumas análises relativas a questões de gênero no Parlamento.\\
Recebeu 4 premiações desde 2012.}%
{A web application that calculate ``similarities'' between political parties based on legislative houses voting data analysis. There are also some gender analysis from the Parliament.
Has been awarded 4 times since 2012.}}\\
\maintechs Django/Python, R, Postgres, nginx, Docker, Gitlab CI, ElasticSearch \eand d3js

\textbf{Transparência Hacker \lang{en}{(\textit{Hacker Transparency})}}\\
\langif{br}%
{Membro da comunidade brasileira Transparência Hacker, que atua nas áreas de Governo Aberto, Transparência Pública, Dados Abertos e afins. A comunidade já criou o Ônibus Hacker, participou da criação do \href{http://blog.openingparliament.org/post/72099651071/a-permanent-hacker-space-in-the-brazilian-congress}{\textbf{Laboratório Hacker do Congresso Nacional}}, influenciou na formulação da Lei de Acesso à Informação e do Marco Civil da Internet, dentre outros.}%
{Member of Brazilian Hacker Transparency community, that advocates for OpenGov, Public Transparency, OpenData, etc. The community created the HackerBus, drove the creation of a \href{http://blog.openingparliament.org/post/72099651071/a-permanent-hacker-space-in-the-brazilian-congress}{\textbf{Congress official HackerLab}}, influenced the Brazilian Public Information Access act and the ``Marco Civil da Internet'', among other projects.}
%
% PUBLICAÇÕES
%
%\sectionlang{br}{Publicações}
%\sectionlang{en}{Publications}
%{\href{http://polignu.org/artigo/twitter-api-elasticsearch-e-kibana-analisando-rede-social}{\textbf{Twitter API, ElasticSearch e Kibana - Analisando a rede social}}\\
%\emph{\langif{br}{Tutorial sobre como coletar, processar e analisar dados do Twitter}{Tutorial on how to gather, process and analyze twitter traffic.}}\\}\\
%%{\href{http://diraol.polignu.org}{\textbf{DiRaOLinux}}\\
%%\emph{\langif{br}{Blog para documentar e compartilhar experiências diárias com software livre}{Blog to document and share daily discovers and issues fixes}}\\}\\
%\href{http://polignu.org/administrativo/caderno-polignu-volume-1-software-e-culturas-livres}{\textbf{Caderno do PoliGNU - Vol 1 - Software e Cultura Livres}}\\
%\emph{\langif{br}{Material introdutório sobre Software Livre e Cultura Livre}{Introductory material about Free Software and Free Culture (pt-br)}}
%
% CURSOS E CERTIFICAÇÕES
%
%\sectionlang{br}{Cursos e Certificações}
%\sectionlang{en}{Courses and Certifications}
%\begin{entrylist}
%  \entry
%    {2015 - 2017}
%    {Data Science Specialization}
%    {\langif{br}{Universidade Johns Hopkins no Coursera}{Johns Hopkins University on Coursera}}
%    {\scriptsize{%
%      \cert{Estatística Inferencial}{Statistical Inference}{https://coursera.org/verify/C2HZUG2HXQJK} \\
%      %
%      \cert{Modelos de Regressão}{Regression Models}{https://coursera.org/verify/T9LTC73MX62A} \\
%      %
%      \cert{Pesquisa reprodutível}{Reproducible Research}{https://coursera.org/verify/U5X4J6YSDC} \\
%      %
%      \cert{Análise Exploratória de Dados}{Exploratory Data Analysis}{https://coursera.org/verify/BCSV3DF2F4} \\
%      %
%      \cert{Obtenção e Limpeza de Dados}{Getting and Cleaning Data}{https://coursera.org/verify/YFY5YWF3MM} \\
%      %
%      \cert{Linguagem R}{R Programming}{https://coursera.org/verify/M2NP642ZUE} \\
%      %
%      \cert{As Ferramentas do Cientista de Dados}{Data Scientist’s Toolbox}{https://coursera.org/verify/ELJEYB2D8A}}} \\
%  \entry
%    {12/2014}
%    {\langif{br}{Objetos Pythonicos}{Pythonic Objects}}
%    {\langif{br}{PythonPro}{PythonPro}}
%    {\href{https://www.python.pro.br/modulos/}{https://www.python.pro.br/modulos/}}
%\end{entrylist}
%
% NOVA PÁGINA
%
\newpage
\pagenumbering{gobble}% Remove page numbers (and reset to 1)
%
% EVENTOS
%
\sectionlang{br}{Eventos}
\sectionlang{en}{Events}

% https://reinvent.awsevents.com/
% https://web.archive.org/web/20181126182100/https://reinvent.awsevents.com/
{\footnotesize{\langif{br}{Nov/18}{Nov/18}}} \- %
\href{https://web.archive.org/web/20181126182100/https://reinvent.awsevents.com/}{%
  \textbf{%
    AWS Re:Invent 2018 @ Las Vegas, \langif{br}{EUA}{USA}%
  }%
}\\
\langif{en}{%
  Attended the Executive Summit invited by AWS Brasil on behalf of EasyTaxi.%
}{%
  Estive no Summit Executivo a convite da AWS Brasil em nome da EasyTaxi.%
}

% https://mastertech.com.br/
% https://web.archive.org/web/20181102114039/https://mastertech.com.br/
{\footnotesize{\langif{br}{Nov/18}{Nov/18}}} \- %
\href{https://mastertech.com.br/}{\textbf{MasterTech}}\\
\langif{en}{%
  Taught a workshop about CI/CD and Containers for a private bank executives.%
}{%
  Ministrei um workshop sobre CI/CD e Containers para executivos de um banco privado.%
}

% http://congressoabraji.blogspot.com/2018/06/ferramentas-facilitam-coleta-de-dados.html
% https://web.archive.org/web/20190504180539/http://congressoabraji.blogspot.com/2018/06/ferramentas-facilitam-coleta-de-dados.html
{\footnotesize{\langif{br}{Jul/18}{Jul/18}}} \- %
\href{http://congressoabraji.blogspot.com/2018/06/ferramentas-facilitam-coleta-de-dados.html}{%
  \textbf{%
    \langif{br}{%
      13\textsuperscript{o} Congresso Internacional de Jornalismo Investigativo%
    }{%
      13\textsuperscript{th} International Congress of Investivative Journalism%
    }%
  }%
}\\
\langif{br}{%
  Apresentação Oficial do projeto Carcerópolis desenvolvido pela ASK-AR para Conectas Direitos Humanos.%
}{%
  Offical announcement of Carcerópolis project, developed by ASK-AR to the Conectas Direitos Humanos NGO.%
}%

% https://sc17.supercomputing.org/
% https://www.amlight.net/?p=2999
% https://web.archive.org/web/20190505183456/https://www.amlight.net/?p=2999
{\footnotesize{\langif{br}{Nov/17}{Nov/17}}} \- \href{https://sc17.supercomputing.org/}{\textbf{The International Conference for High Performance Computing, Networking, Storage and Analysis - SC17 @ Denver, \langif{br}{EUA}{USA}}}\\
\langif{br}{%
  Apresentação do projeto Kytos desenvolvido na UNESP, com execução de uma
  \href{https://web.archive.org/web/20190505183456/https://www.amlight.net/?p=2999}{demonstração}
  de transmissão de dados entre São Paulo e Miami atingindo 370Gpbs numa
  parceria entre Caltech, AmLight, RNP e UNESP%
}{%
  Attended the Conference to present Kytos.
  During the conference we
  \href{https://web.archive.org/web/20190505183456/https://www.amlight.net/?p=2999}{demoed}
  a data transmission between South (São Paulo) and North (Miami) hemispheres
  reaching a steady rate of 370Gbps, in a joint effort between Caltech,
  AmLight, RNP and UNESP.%
}%

% https://web.archive.org/web/20190203153054/https://tnc17.geant.org/
{\footnotesize{\langif{br}{Jun/17}{Jun/17}}} \- \href{https://tnc17.geant.org/}{\textbf{TNC17 Networking Conference 2017 @ Linz, Austria}}\\
\langif{br}{%
  Apresentação do projeto Kytos desenvolvido na UNESP.%
}{%
  Attended the Conference to present Kytos (Software Defined Controller developed at UNESP).%
}%

% {\footnotesize{\langif{br}{Dez/16}{Dec/16}}} - \href{https://www.facebook.com/events/1617891821849460/}{\textbf{UP\{Design For All\}it}}\\
% \langif{br}{Mentor do 3$^o$ encontro \textit{UPWIT - Design For All}. Um evento com objetivo de pensar tecnologias de forma inclusiva}
%            {Mentor on the 3$^{rd}$ \textit{UPWIT - Design For All}. An event to think inclusive technologies}.
%
{\footnotesize{\langif{br}{Mai/16}{May/16}}} - \href{https://rgaiacs.github.io/2016-05-27-ccsl/}{\textbf{Software Carpentry Workshop}}\\
\langif{br}{Instrutor de R no workshop realizado na Universidade de São Paulo (USP)}
           {R instructor on the workshop held at University of São Paulo}.
% http://web.archive.org/web/20170301105603/https://rgaiacs.github.io/2016-05-27-ccsl/

{\footnotesize{\langif{br}{Dez/15}{Dec/15}}} - \href{http://www.lets.go.gov.br/}{\textbf{Hackatona Let's GO}}\\
\langif{br}{Mentor na Hackathona promovida pela Secretaria de Estado de Gestão e Planejamento do Governo de Goiás}
           {Mentor on the Hackathon promoted by Goiás State Government Planning and Management Secretary}.
% http://web.archive.org/web/20160504011426/http://www.lets.go.gov.br/

{\footnotesize{\langif{br}{Out/15}{Oct/15}}} - \href{http://www.hackagenda.com.br/evento/sesi-cultura-digital-2015-hackathon-o-museu-do-seculo-xxi/}{\textbf{SESI Cultura Digital 2015 + Hackathon “O museu do século XXI”}}\\
\langif{br}{Instrutor do workshop ``mapeamento com \textbf{CartoDB}'' e mentor dos grupos participantes do Hackathon}
           {``Mapping with CartoDB'' workshop instructor and Hackathon mentor}.
% http://web.archive.org/web/20170301110629/http://pinghacker.com.br/hackagenda/sesi-cultura-digital-2015-hackathon-o-museu-do-seculo-xxi/
% http://web.archive.org/web/20170301110501/http://www.firjan.com.br/sesi/qualidade-de-vida/projetos-sesi-cultural/sesi-cultura-digital/hackaton/

% {\footnotesize{\langif{br}{Set/15}{Sep/15}}} - \href{http://rgaiacs.github.io/2015-09-10-usp/}{\textbf{Software Carpentry Workshop}}\\
% \langif{br}{Instrutor de git no workshop realizado na Universidade de São Paulo (USP)}
%            {Git instructor on the workshop held at University of São Paulo}.
% http://web.archive.org/web/20160602202812/https://rgaiacs.github.io/2015-09-10-usp/

{\footnotesize{\langif{br}{Ago/15}{Aug/15}}} - \href{https://portaldoaluno.casperlibero.edu.br/agenda-eventos/semana-de-jornalismo-casper-libero-2/}{\textbf{\langif{br}{Seminário de Jornalismo da Faculdade Casper Líbero}{Journalism Seminar on Casper Líbero University}}}\\
\langif{br}{Debatedor na mesa Jornalismo de Dados e Políticas Públicas de Comunicação}
           {Debater in a panel on Data Journalism and Public Policies for Communications}.
% https://web.archive.org/web/20170301110810/https://portaldoaluno.casperlibero.edu.br/agenda-eventos/semana-de-jornalismo-casper-libero-2/

{\footnotesize{Jun/15}} - \href{https://setemasters.imasters.com.br/edicoes/design-de-apis/}{\langif{br}{\textbf{7Masters Design de APIs}}{\textbf{7Masters APIs Design}}}\\
\langif{br}{Apresentação: APIs para visualização de dados públicos}%
           {Talk: APIs for open government data visualizations}.
% https://setemasters.imasters.com.br/edicoes/design-de-apis/
% https://web.archive.org/web/20160624032825/http://setemasters.imasters.com.br/edicoes/design-de-apis/
% https://setemasters.imasters.com.br/conversas/apis-para-visualizacao-de-dados-publicos/
% https://web.archive.org/web/20170301120904/https://setemasters.imasters.com.br/conversas/apis-para-visualizacao-de-dados-publicos/

{\footnotesize{Jul/14}} - \href{http://abraji.org.br/congresso/}{\textbf{\langif{br}{9$^o$ Congresso Internacional de Jornalismo Investigativo}{9$^{th}$ International Investigative Journalism Congress}}}\\
\langif{br}{Instrutor do workshop ``Programação para Jornalistas''}
           {``Coding for journalists'' workshop instructor}.

{\footnotesize{Jan/14}} - \href{http://www.techtudo.com.br/noticias/noticia/2014/01/cp2014-campus-party-brasil-2014-guia-traz-principais-atividades-por-perfil.html}{\textbf{\langif{br}{Campus Party Brasil 2014}{2014 Brazil Campus Party}}}\\
\langif{br}{Oficina: ``Dinheiro Público, para onde vai, para onde você quer que vá?''}
           {Workshop: ``Public Budget, where does it go, where do you want it to go to?''}
% https://web.archive.org/web/20150517063722/http://www.techtudo.com.br/noticias/noticia/2014/01/cp2014-campus-party-brasil-2014-guia-traz-principais-atividades-por-perfil.html

{\footnotesize{Nov/13}} - \href{https://knightcenter.utexas.edu/pt-br/blog/00-14390-inscricoes-abertas-para-primeiro-curso-da-anj-com-o-centro-knight-introducao-ao-jornal}{\textbf{\langif{br}{Curso de Introdução ao Jornalismo de Dados}{Introduction to Data Journalism MOOC}}}\\
\langif{br}{Instrutor do módulo ``Programação''. Curso organizado pelo Knight Center e pela \nohyphens{Associação} Brasileira de Jornais}
           {``How to code'' module instructor. Course organized by Knight Center and the Brazilian National Journals Association}.
%https://web.archive.org/web/20170301112535/https://knightcenter.utexas.edu/pt-br/blog/00-14390-inscricoes-abertas-para-primeiro-curso-da-anj-com-o-centro-knight-introducao-ao-jornal

{\footnotesize{Nov/13}} - \href{https://web.archive.org/web/20140330181037/http://2.encontro.dados.gov.br/encontro.html}{\textbf{\langif{br}{II Encontro Nacional de Dados Abertos}{II Brazilian National Conference of OpenData}}}\\
\langif{br}{\textit{Chairperson} da Trilha de Visualização de Dados}
           {Dataviz trail Chairperson}.
% http://2.encontro.dados.gov.br/encontro.html

{\footnotesize{Jul/13}} - \href{http://rodadahacker.com/}{\textbf{II RodAda Hacker}}\\
\langif{br}{Mentor na segunda edição da ``RodAda Hacker, um workshop de um dia que objetiva promover autonomia  tecnológica a mulheres}%
           {Mentor on the second edition of ``RodAda Hacker'', a one day workshop that aims to promote women technological autonomy}.

{\footnotesize{\langif{br}{Mai/13}{May/13}}} - \href{http://www2013.org/}{\textbf{www2013 Conference}}\\
\langif{br}{Palestrante na mesa ``\textit{eGov and Open Data Camp}'', falando sobre ``\textit{The challenges and opportunities of data journalism}''}
           {Speaker on the ``eGov and Open Data Camp'' panel, talking about ``The challenges and opportunities of data journalism''}.
% https://web.archive.org/web/20161223084334/http://www2013.org/tracks/w3c-track/

% {\footnotesize{\langif{br}{Abr/13}{Apr/13}}} - \href{https://setemasters.imasters.com.br/edicoes/opendata/}{\textbf{7Masters OpenData}}\\
% \langif{br}{Apresentação do projeto Radar Parlamentar.}
%            {Speech about the Parliamentary Radar Project}.
% https://setemasters.imasters.com.br/edicoes/opendata/
% https://web.archive.org/web/20140923074920/http://setemasters.imasters.com.br/edicoes/opendata/
% https://setemasters.imasters.com.br/conversas/radar-parlamentar/
% https://web.archive.org/web/20170301120109/https://setemasters.imasters.com.br/conversas/radar-parlamentar/

{\footnotesize{Mar/13}} - \href{http://rodadahacker.com/}{\textbf{I RodAda Hacker}}\\
\langif{br}{Mentor na primeira edição da ``RodAda Hacker, um workshop de um dia que objetiva dar a mulheres mais autonomia com tecnologias}%
           {Mentor on the first edition of ``RodAda Hacker'', a one day workshop that aims to give woman autonomy with technologies}.

% {\footnotesize{Nov/12}} - \href{http://hackshackers.com/blog/author/gfaleiros/}{\textbf{\langif{br}{1$^o$ Hackathon HacksHackers}{1$^{st}$ HacksHackers Hackathon}}}
% https://web.archive.org/web/20160414124730/http://hackshackers.com/blog/author/gfaleiros/

% {\footnotesize{\langif{br}{Out/12}{Oct/12}}} - \href{http://yapcbrasil.org.br/2012/talk/110}{\textbf{YAPC::Brasil 2012}}\\
% \langif{br}{Palestrante: ``Radar Parlamentar: Desvendando a política legislativa''}%
%            {Speaker: ``The Parliamentary Radar: Unraveling the legislative politics''}.
% https://web.archive.org/web/20121022134231/http://yapcbrasil.org.br/2012/talks

% {\footnotesize{Jul/12}} - \href{http://softwarelivre.org/fisl13}{\textbf{\langif{br}{13$^o$ Fórum Internacional de Software Livre (FISL)}{13$^{th}$ International Free Software Forum}}}\\
% \langif{br}{Apresentação do Projeto Radar Parlamentar no Workshop de Software Livre (WSL)}
%            {Presented the Parliamentary Radar project on the Free Software Workshop (WSL)}.
% http://softwarelivre.org/articles/0125/7499/caderno-digital.pdf

{\footnotesize{\langif{br}{Mai/10}{May/10}}} - \href{http://artigo19.org/infoedireitoseu/?p=560}{\textbf{\langif{br}{O direito ao Acesso à Informação}{The Right to Information Access}}}\\
\langif{br}{Palestrante sobre o Direito ao Acesso à Informação num debate na PUC-SP}%
           {Speaker on the ``Right to Information Access'' debate at PUC University}.
% https://web.archive.org/web/20170120024235/http://artigo19.org/infoedireitoseu/?p=560

%
%\\
%
%\begin{flushleft}
%\emph{November 22th, 2014}
%\end{flushleft}
%\begin{flushright}
%
%\end{flushright}
%
%%% This piece of code has been commented by Karol Kozioł due to biblatex errors.
%
%\printbibsection{article}{article in peer-reviewed journal}
%\begin{refsection}
%  \nocite{*}
%  \printbibliography[sorting=chronological, type=inproceedings, title={international peer-reviewed conferences/proceedings}, notkeyword={france}, heading=subbibliography]
%\end{refsection}
%\begin{refsection}
%  \nocite{*}
%  \printbibliography[sorting=chronological, type=inproceedings, title={local peer-reviewed conferences/proceedings}, keyword={france}, heading=subbibliography]
%\end{refsection}
%\printbibsection{misc}{other publications}
%\printbibsection{report}{research reports}

\end{document}
